\chapter{Fase 0: el plan, las reglas y la computadora}
\label{cap:d-f0}
\textbf{Fechas}: 26 y 27 de septiembre de 2026. \textbf{Notas}: \codigo{00-fundacion.md}, \codigo{02-entorno.md}.

\textbf{Qué se buscaba}: decidir qué proyecto hacer, con qué reglas y cómo trabajar, y dejar la computadora lista para
procesar millones de datos con Spark.

\begin{opciones}[0.1 · El estado · 26-sep-2026]
\textbf{La pregunta}: el Problema Prototípico pide elegir un estado de México.

\textbf{Lo que eligió Brandon}: \textbf{Quintana Roo}, y su nombre como autor en todo el código (``Por ahora ponme a mi
Brandon Uriel Garcia Sanchez'').

\textbf{Consecuencia}: el estado más visitado del Caribe mexicano, con un problema de saturación visible en 2026 (sargazo,
Tulum, Cozumel) y con fuentes oficiales abundantes (SITUR-Q, DataTur, INAH).
\end{opciones}

\begin{opciones}[0.2 · Empezar desde cero · 26-sep-2026]
\textbf{La pregunta}: ya existía un proyecto anterior (el paquete \codigo{cauce}, con modelos, personas y cifras). ¿Se
reutiliza?

\textbf{Lo que eligió Brandon}: \emph{``Quiero que armemos desde 0 todos los sistemas que consultaste esos no me gustaron y
quiero que tu lo hagas otra vez desde 0 y lo armemos los dos juntos''}.

\textbf{Descartado}: portar el paquete anterior y reentrenar sus modelos.

\textbf{Consecuencia}: ninguna cifra, modelo ni flujo del proyecto anterior se usa. Cada técnica se diseñó y validó de
nuevo, fase por fase. La carpeta antigua quedó solo como referencia histórica y excluida del mapa del proyecto.
\end{opciones}

\begin{opciones}[0.3 · Qué sistema construir: la fusión de tres alternativas · 27-sep-2026]
\textbf{La pregunta}: Brandon pidió \emph{``5 ALTERNATIVAS CON CADA RUBRICA CUMPLIDA, DONDE YA SABES QUE ESTAN LOS DATOS''},
para no descubrir a medio camino que una no era posible. Se presentaron cinco, cada una con sus datos verificados:

\begin{center}
\small
\begin{tabularx}{\linewidth}{lY}
\encabezado \h{Alternativa} & \h{Idea central} \\
A1 Radar & Medir dónde hay presión y dónde hay espacio hoy \\
\filagris A2 Voz del viajero & Construir la campaña desde las reseñas de los turistas \\
A3 Pronóstico & Predecir cuándo conviene ir y cuánto invertir \\
\filagris A4 Portafolio & Tratar los destinos como un portafolio de inversión \\
A5 Tiempo casi real & Ajustar la campaña semana a semana mientras corre \\
\bottomrule
\end{tabularx}
\normalsize
\end{center}

\textbf{Lo que eligió Brandon}: \emph{``A5 A1 y A3 como podemos fucionarlo dame el plan donde no se encimen ambas y se
compaginen para no estar saturado y funcionen en la pagina''}.

\textbf{Cómo se evitó que se enciman}: cada módulo responde una sola pregunta en un solo horizonte de tiempo: el Radar,
\emph{dónde} hoy (semanas); el Pronóstico, \emph{cuándo} y \emph{cuánto} (1 a 12 meses); la Torre, \emph{qué hacer esta
semana}. Markov trabaja solo en semanas y Monte Carlo solo en meses; el clasificador predice un estado y el pronóstico un
número; hay un solo modelo de optimización con dos etapas, y la Torre no entrena modelos propios.

\textbf{Descartado}: A2 y A4 como ejes. La minería de texto de A2 se conservó para construir a las viajeras ideales.

\textbf{Evidencia de que había datos}: DataTur publica 135 archivos semanales de ocupación con 7 centros de Quintana Roo;
SITUR-Q responde 45 indicadores por su API; la NOAA tiene las tormentas desde 1851.
\end{opciones}

\begin{opciones}[0.4 · Cómo trabajar · 26-sep-2026]
\textbf{Lo que pidió Brandon}: \emph{``necesitamos que lo desarrolles CONMIGO Y NO VIBEDEES TODO SIN SABER COMO CAJA
NEGRO para saber que haces y el porque''} y \emph{``me expliques el porque tomaste una decisión y el porque lo hiciste
argumentandomelo con datos reales y no supociciones o alucinaciones''}.

\textbf{Consecuencia}: el protocolo de cinco pasos (cifra real, opciones, Brandon elige, pieza pequeña, nota escrita), una
fase a la vez, y las reglas de oro: no inventar datos, toda cifra rastreable, lo estimado se marca con \codigo{_est}, nada
de \emph{scraping} prohibido, y si falta un dato, se detiene y se avisa. Todo quedó en \codigo{OBJETIVO.md} y
\codigo{CLAUDE.md}.
\end{opciones}

\begin{opciones}[0.5 · Ecuaciones y ``cómo lo resolví'' · 27-sep-2026]
\textbf{Lo que pidió Brandon}: \emph{``OTra cosa que me falto decirte es que tengas las ecuaciones y como lo resolviste
porque es un proyecto escolar sabes :)''}.

\textbf{Consecuencia}: cada modelo se documenta con cinco partes: la ecuación, sus supuestos, el método de resolución paso
a paso, un ejemplo resuelto a mano con números reales y dónde está en el código. Todo en
\codigo{docs/metodologia/ECUACIONES.md}, en los notebooks y en la \emph{Guía técnica}.
\end{opciones}

\begin{opciones}[0.6 · Herramientas de trabajo · 27-sep-2026]
\textbf{Lo que eligió Brandon}: repositorio \textbf{git} (historial de cada cambio) y \textbf{Graphify}, un mapa de
conocimiento que conecta reglas, fuentes, módulos y regiones. La carpeta del proyecto anterior queda \textbf{fuera} del
mapa.

\textbf{Más tarde} (30-sep-2026): el repositorio se hizo \textbf{público} y la página se publicó en GitHub Pages
(decisión de Brandon). Los datos crudos nunca se suben.
\end{opciones}

\begin{opciones}[0.7 · La computadora: Spark en Windows · 27-sep-2026]
\textbf{La pregunta}: la máquina solo tenía Java 1.8. PySpark 4, lo que instala pip por defecto, exige Java 17. Y en Windows,
Spark no escribe Parquet sin dos archivos de Hadoop (\codigo{winutils.exe} y \codigo{hadoop.dll}).

\begin{center}
\small
\begin{tabularx}{\linewidth}{>{\raggedright\arraybackslash}p{0.3\linewidth}YY}
\encabezado \h{Opción} & \h{A favor} & \h{En contra} \\
\textbf{Entorno propio: PySpark 3.5.6 + Java 17 + Hadoop 3.3.6 (elegida)} & Versiones compatibles y verificadas; Java 17
solo dentro del proyecto & Hay que instalar Java 17 \\
\filagris Spark dentro de Docker & Funciona igual en cualquier computadora & Pesado de arrancar y de explicar el día del
coloquio \\
Cambiar el Java de todo Windows & Simple & Podía romper otros programas que usan Java 8 \\
\bottomrule
\end{tabularx}
\normalsize
\end{center}

\textbf{Evidencia}: la prueba de humo (\codigo{tests/test_entorno.py}) leyó un CSV de 3 filas con Spark, escribió Parquet,
lo volvió a leer y la suma dio $1+2+3=6$, la esperada. Los dos archivos de Hadoop se guardaron con su huella SHA-256.
\end{opciones}

\begin{opciones}[0.8 · No promover destinos de playa en 2026 · 27-sep-2026]
\textbf{La pregunta}: ¿qué tipo de destino promover?

\textbf{Evidencia}: sargazo récord (más de 104,700 toneladas; 56 de 140 playas en rojo en agosto, y la franja más
afectada va de Tulum a Xcalak); crisis de Tulum (ventas hasta $-60\,\%$, cierres); Cozumel declarado lleno por cruceros;
Isla Mujeres con carencias de agua y drenaje; Holbox con más de 90,000 toneladas de basura acumulada.

\textbf{Consecuencia}: la campaña deja de ser de playa. Se vuelve una ruta de cultura, bahía, laguna y comunidad en el
sur, que se confirmaría con una tabla de criterios calculada con datos en la Fase 3.
\end{opciones}

\begin{opciones}[0.9 · Un documento ejecutivo para el equipo · 28-sep-2026]
\textbf{La decisión}: al cerrar cada fase se agrega un capítulo en lenguaje no técnico, con el estilo de la UNRC, para que
otra integrante del equipo pueda redactar el entregable en limpio. Se genera como PDF con un comando.
\end{opciones}
